\chapter{Diseño del sistema}

\begin{quotation}[Book]{A guide to the Project Management Body of Knowledge (PMBOK guide) (2017)}

\end{quotation}

\begin{abstract}
Se comentará el diseño del sistema desde el punto de vista de los actores implicados así como de los modelos que nutriran la implementación a posteriori del sistema.
\end{abstract}

\clearpage

\section{Modelos diseñados}

A lo largo del ciclo de vida del proyecto nos hemos encontrado ante numerosos cambios en el sistema que se ha implemetnado para la Gestión de las Permutas, al encontrarnos ante un caso real de un producto real que pasará a producción mantener numerosas reuniones al principio del ciclo de vida para detectar las posibles carencias así como las modificaciones pertinentes en los requisitos y modelos nos ahorrarían cambios que a futuro serían más costosos. 

Es por ello que se presentan los modelos generados con cada cambio así como las modificaciones que se han realizado y el sentido de estas en cada etapa del proyecto.

\textbf{Versión original}


El primer diseño del sistema tienen un único objetivo, modelar el sistema de permutas como se tenía pensado por parte de los desarrolladores.\ref{fig:diseno00}.

\begin{figure}[htbp]
\begin{center}

\caption{Diagrama UML de diseño original}
\label{fig:diseno00}
\end{center}
\end{figure}

\textbf{Primera revisión}


Tras la segunda reunión celebrada el día 8 de enero y en aras de reflejar el sistema de notificaciones e incidencias se amplia y modifica el diseño del sistema para que refleje la nueva realidad del proyecto.\ref{fig:diseno00}.

\begin{figure}[htbp]
\begin{center}

\caption{Diagrama UML de diseño tras la primera revisión}
\label{fig:diseno00}
\end{center}
\end{figure}


\section{Diagrama de actores}

Para el diseño del diagrama de actores se han tenido en cuenta los siguientes requisitos:
\begin{itemize}
\item El Estudiante puede Buscar y Solicitar permutas, así como Aceptar o Rechazar las que se le propongan y Consultar el estado de las solicitudes que haya hecho.
\item La Secretaría es la autoridad encargada de Validar o Rechazar formalmente la permuta una vez que los estudiantes estén de acuerdo.
\item El Servicio de Mantenimiento atiende las incidencias relacionadas con el sistema para que el servicio de permutas funcione correctamente.
\end{itemize}

Este es por tanto el resultado del diagrama \ref{fig:diseno00}.

\begin{figure}[htbp]
\begin{center}

\caption{Diagrama de actores tras original}
\label{fig:diseno00}
\end{center}
\end{figure}